% Physical unit 025; accepted Session 03 after fresh 400-dpi authority replay.
\GordanUnitStart{025}{printed page 20}
\GordanTopFolio{20}
\begingroup\small
\[
\begin{array}{ll}
\text{either } l<l_1, & \psi>\omega,\\
\text{or } l>l_1, & \psi>\omega,\\
\text{or } l<l_1, & \psi<\omega,\\
\text{or } l>l_1, & \psi<\omega.
\end{array}
\]
The equations of the lines of the first and fourth genus take this form:
\[
 \psi-\omega=\int_{\varphi_1}^{\varphi}\cos l_0\,d\varphi\,{\sqrt{1-e^2+e^2\sin^2l_0\cos^2\varphi}\over 1-\sin^2l_0\cos^2\varphi}
\]
and of the second and third this one:
\[
 \omega-\psi=\int_{\varphi_1}^{\varphi}\cos l_0\,d\varphi\,{\sqrt{1-e^2+e^2\sin^2l_0\cos^2\varphi}\over 1-\sin^2l_0\cos^2\varphi}.
\]

\begin{center}\S\ 3.\end{center}

The nature of all geodetic lines is this: each one contains all the points of the zone of the spheroid enclosed by the circles $l_0$ and $-l_0$.

\begin{center}ARGUMENT.\end{center}

It must be shown that any geodetic line whatever cuts every circle of the zone successively at all its points; hence that their longitudes successively assume all numbers lying between $0$ and $2\Pi$.

The longitudes of this section are sought from the equation
\[
 \psi-\omega=\pm\int_{\varphi_1}^{\varphi}\cos l_0\,d\varphi\,{\sqrt{1-e^2+e^2\sin^2l_0\cos^2\varphi}\over 1-\sin^2l_0\cos^2\varphi}
\]
in which $\varphi_1$ is constant, while the values of the angle $\varphi$ are derived from the equation
\[
 \sin l=\sin l_0\cos\varphi.
\]
Putting, in the first section,
\[
 \varphi=\varphi_2\qquad \psi=\psi_2,
\]
in the following sections successively the angle becomes
\[
 \varphi=\varphi_2+2\pi\qquad \varphi_2+4\pi
\]
\[
 \psi=\omega+\int_{\varphi_1}^{\varphi_2+2\pi}\cos l_0\,d\varphi\,{\sqrt{1-e^2+e^2\sin^2l_0\cos^2\varphi}\over 1-\sin^2l_0\cos^2\varphi}
\]
\[
 \omega+\int_{\varphi_1}^{\varphi_2+4\pi}\cos l_0\,d\varphi\,{\sqrt{1-e^2+e^2\sin^2l_0\cos^2\varphi}\over 1-\sin^2l_0\cos^2\varphi}
\]
or $\omega+4\mathcal A(l_0)$ \qquad $\omega+8\mathcal A(l_0)$,
\endgroup
